\bookchapter{附錄C：極簡術語“黑話”對照表}

\begin{center}
\begin{tabularx}{\linewidth}{|X|X|X|}
\hline
\textbf{君或聽過之詞} & \textbf{實則即指……} & \textbf{一句話解釋} \tabularnewline
\hline
SoC / 芯片 & 手機之大腦與心臟 & 決定手機快否、遊戲卡否之核心部件。 \tabularnewline
RAM (運存) & 手機之"短期工作桌" & 桌越大，能同時開之App越多，切換越流暢。 \tabularnewline
ROM (存儲) & 手機之"大倉庫" & 存照片、視頻、App之處。128G、256G言之即它。 \tabularnewline
OLED / LCD & 屏幕之兩種材質 & OLED色彩更豔、更省電；LCD更護眼、成本低。 \tabularnewline
NFC & 近距離無線通訊 & 令君手機能化身公交卡、門禁卡、銀行卡之"貼貼"功能。 \tabularnewline
OTA更新 & 無線系統升級 & 於【設置】裏點之"系統更新"，即它。 \tabularnewline
Root & 安卓超級管理員權限 & 得之後可改系統文件，然亦有風險。 \tabularnewline
BL鎖 & Bootloader鎖 & 廠商加之安全鎖，解鎖後方能刷機。 \tabularnewline
NPU & 神經網絡處理器 & 專門處理AI任務之芯片，較CPU快數十倍、功耗更低。 \tabularnewline
端側AI & 手機本地運行之AI & AI於手機上徑行運行，不需聯網，速快、隱私好。 \tabularnewline
Qi2 & 新一代無線充電標準 & 帶磁吸定位之無線充電，類似MagSafe，充電更便。 \tabularnewline
eSIM & 嵌入式SIM卡 & 無實體卡之SIM，徑寫入手機，換運營商不用換卡。 \tabularnewline
LTPO & 自適應刷新率技術 & 屏幕刷新率依內容自動調節（1-120Hz），流暢又省電。 \tabularnewline
\hline
\end{tabularx}
\end{center}
